\subsubsection{General Mathematical Derivation}\label{sec:bernstein-vazirani-general-mathematical-derivation}

The previous sections explained the Bernstein-Vazirani algorithm intuitively. We now prove mathematically that, for any hidden string:
\begin{equation*}
    w \in \left\{0,1\right\}^n
\end{equation*}
The final state of the data register (i.e. the first $n$ qubits) is exactly \ket{w}. The derivation is based on the $n$-qubit Hadamard identity:
\begin{equation}
    H^{\otimes n} \ket{x} = \frac{1}{\sqrt{2^n}} \sum_{y \in \left\{0,1\right\}^n} (-1)^{x \cdot y} \ket{y}
\end{equation}
Where:
\begin{itemize}
    \item $x \in \left\{0,1\right\}^n$ is an $n$-bit string.
    \item $y \in \left\{0,1\right\}^n$ is an $n$-bit string.
    \item $x \cdot y = x_1 y_1 \oplus \dots \oplus x_n y_n$ is the bitwise inner product of $x$ and $y$, modulo 2.
    \item $(-1)^{x \cdot y}$ is a phase factor that is either $+1$ or $-1$, depending on the value of $x \cdot y$.
\end{itemize}

\highspace
\begin{flushleft}
    \textcolor{Green3}{\faIcon{play-circle} \textbf{Initial State}}
\end{flushleft}
The algorithm starts from:
\begin{equation}
    \ket{\psi_0} = \ket{0}^{\otimes n} \ket{1}
\end{equation}
The first $n$ qubits form the \textbf{data register}, while the final qubit is the ancilla.

\highspace
\begin{flushleft}
    \textcolor{Green3}{\faIcon{project-diagram} \textbf{After the First Hadamard Layer}}
\end{flushleft}
We apply:
\begin{equation*}
    H^{\otimes n} H
\end{equation*}
Where the first $n$ Hadamards act on the data register, and the final Hadamard acts on the ancilla. Thus, for the data register we have:
\begin{equation*}
    H^{\otimes n} \ket{0}^{\otimes n} = \frac{1}{\sqrt{2^n}} \sum_{x \in \left\{0,1\right\}^n} (-1)^{0 \cdot x} \ket{x} = \frac{1}{\sqrt{2^n}} \sum_{x \in \left\{0,1\right\}^n} \ket{x}
\end{equation*}
Instead, for the ancilla we have:
\begin{equation*}
    H \ket{1} = \frac{1}{\sqrt{2}} \left( \ket{0} - \ket{1} \right) = \ket{-}
\end{equation*}
Thus, the state after the first Hadamard layer is:
\begin{equation}
    \ket{\psi_1} = \frac{1}{\sqrt{2^n}} \sum_{x \in \left\{0,1\right\}^n} \ket{x} \ket{-}
\end{equation}
The \hl{data register} is now in a \hl{uniform superposition of all possible input strings}, while the ancilla is in the state \ket{-}, which is prepared for phase kickback.

\highspace
\begin{flushleft}
    \textcolor{Green3}{\faIcon{magic} \textbf{After the Oracle}}
\end{flushleft}
The oracle satisfies:
\begin{equation*}
    U_f \ket{x} \ket{y} = \ket{x} \ket{y \oplus f(x)}
\end{equation*}
Since the ancilla is in the state \ket{-}, \textbf{phase kickback} gives (complete derivation on \autopageref{eq:bernstein-vazirani-oracle-action}):
\begin{equation*}
    U_f \ket{x} \ket{-} = (-1)^{f(x) = w \cdot x} \ket{x} \ket{-}
\end{equation*}
By linearity, the complete state after the oracle is:
\begin{equation}
    \ket{\psi_2} = \frac{1}{\sqrt{2^n}} \sum_{x \in \left\{0,1\right\}^n} (-1)^{w \cdot x} \ket{x} \ket{-}
\end{equation}
At this point, the \hl{hidden string is encoded entirely in the phases of the data register.}

\highspace
\begin{flushleft}
    \textcolor{Green3}{\faIcon{compress-arrows-alt} \textbf{Applying the Final Hadamard Layer}}
\end{flushleft}
We now apply $H^{\otimes n}$ to the data register:
\begin{equation*}
    \ket{\psi_3} = \frac{1}{\sqrt{2^n}} \sum_{x \in \left\{0,1\right\}^n} (-1)^{w \cdot x} H^{\otimes n} \ket{x} \ket{-}
\end{equation*}
Using the parallel Hadamard identity:
\begin{equation*}
    H^{\otimes n} \ket{x} = \frac{1}{\sqrt{2^n}} \sum_{y \in \left\{0,1\right\}^n} (-1)^{x \cdot y} \ket{y}
\end{equation*}
We obtain:
\begin{equation*}
    \ket{\psi_3} = \frac{1}{2^n} \sum_{x \in \left\{0,1\right\}^n} \sum_{y \in \left\{0,1\right\}^n} (-1)^{w \cdot x} (-1)^{x \cdot y} \ket{y} \ket{-}
\end{equation*}
Rearranging the sums:
\begin{equation*}
    \ket{\psi_3} = \frac{1}{2^n} \sum_{y \in \left\{0,1\right\}^n} \left( \sum_{x \in \left\{0,1\right\}^n} (-1)^{w \cdot x} (-1)^{x \cdot y} \right) \ket{y} \ket{-}
\end{equation*}
Thus, the amplitude of each basis state \ket{y} is:
\begin{equation*}
    \frac{1}{2^n} \sum_{x \in \left\{0,1\right\}^n} (-1)^{w \cdot x} (-1)^{x \cdot y}
\end{equation*}
Because the exponents are computed modulo 2:
\begin{equation*}
    \left(-1\right)^{a} \left(-1\right)^{b} = \left(-1\right)^{a \oplus b}
\end{equation*}
Therefore, we can combine the exponents:
\begin{equation*}
    \left(-1\right)^{w \cdot x} \left(-1\right)^{x \cdot y} = \left(-1\right)^{w \cdot x \oplus x \cdot y}
\end{equation*}
The scalar product modulo 2 is linear, so we can factor out $x$:
\begin{equation*}
    w \cdot x \oplus x \cdot y = x \cdot (w \oplus y)
\end{equation*}
Thus, the state after the final Hadamard layer is:
\begin{equation}
    \ket{\psi_3} = \frac{1}{2^n} \sum_{y \in \left\{0,1\right\}^n} \left( \sum_{x \in \left\{0,1\right\}^n} (-1)^{x \cdot (w \oplus y)} \right) \ket{y} \ket{-}
\end{equation}
Everything now depends on the inner sum:
\begin{equation*}
    \sum_{x \in \left\{0,1\right\}^n} (-1)^{x \cdot (w \oplus y)}
\end{equation*}

\begin{flushleft}
    \textcolor{Green3}{\faIcon{check-circle} \textbf{Measurement}}
\end{flushleft}
At this point, there are two cases:
\begin{itemize}
    \item If $w = y$, then $w \oplus y = 0^n$, therefore the inner sum is:
    \begin{equation*}
        \sum_{x \in \left\{0,1\right\}^n} (-1)^{x \cdot 0^n} = \sum_{x \in \left\{0,1\right\}^n} 1 = 2^n
    \end{equation*}
    So it becomes a giant sum of $1$ terms, until it reaches $2^n$. The amplitude of \ket{w} is therefore:
    \begin{equation*}
        \frac{1}{2^n} \cdot 2^n = 1
    \end{equation*}
    \begin{deepeningbox}[: How did we measure the amplitude?]
        After the final Hadamard layer we have:
        \begin{equation*}
            \ket{\psi_3} = \frac{1}{2^n} \sum_{y \in \left\{0,1\right\}^n} \left( \sum_{x \in \left\{0,1\right\}^n} (-1)^{x \cdot (w \oplus y)} \right) \ket{y} \ket{-}
        \end{equation*}
        This equation shows an interesting structure:
        \begin{align*}
            \ket{\psi_3} &= A (00 \dots 0)\ket{00 \dots 0} + \\
            & \phantom{=} \,\,\, A (00 \dots 1)\ket{00 \dots 1} + \\
            & \phantom{=} \,\,\, \dots + A (11 \dots 1)\ket{11 \dots 1}
        \end{align*}
        Where $A(y)$ is the amplitude of the basis state \ket{y}:
        \begin{equation*}
            A(y) = \frac{1}{2^n} \sum_{x \in \left\{0,1\right\}^n} (-1)^{x \cdot (w \oplus y)}
        \end{equation*}
        So \textbf{each basis state has its own amplitude}. Now, \emph{what is the amplitude of the basis state \ket{y}?} There are $2^n$ possibilities, and instead of writing all of them, we compute them generically. If $y = w$, then $w \oplus y = 0^n$, and the amplitude of \ket{w} is:
        \begin{equation*}
            A(w) = \frac{1}{2^n} \sum_{x \in \left\{0,1\right\}^n} (-1)^{x \cdot 0^n} = \frac{1}{2^n} \sum_{x \in \left\{0,1\right\}^n} 1 = \frac{1}{2^n} \cdot 2^n = 1
        \end{equation*}
    \end{deepeningbox}


    \item If $y \neq w$, then $w \oplus y \neq 0^n$, because at least one bit of $w \oplus y$ is equal to $1$. As $x$ ranges over all $2^n$ possible $n$-bitstrings, exactly half of the values satisfy $x \cdot (w \oplus y) = 0$ and the other half satisfy $x \cdot (w \oplus y) = 1$. Therefore, the inner sum contains equally many $+1$ and $-1$ terms:
    \begin{equation*}
        \sum_{x \in \left\{0,1\right\}^n} (-1)^{x \cdot (w \oplus y)} = 2^{n-1} - 2^{n-1} = 0
    \end{equation*}
    The amplitude of every state \ket{y} with $y \neq w$ \textbf{vanishes}. This is the mathematical expression of \textbf{destructive interference}.
\end{itemize}
Only the term with $y = w$ has nonzero amplitudes. Therefore:
\begin{equation*}
    \ket{\psi_3} = \frac{1}{2^n} \left(2^n \ket{w}\right)\ket{-} = \ket{w} \ket{-}
\end{equation*}
Thus:
\begin{equation}
    \ket{\psi_{\text{final}}} = \ket{w} \ket{-}
\end{equation}
The data register is exactly in the computational basis state corresponding to the hidden bitstring. Consequently, \hl{measuring the data register will yield the hidden string $w$ with probability $1$.}

\highspace
\begin{takeawaysbox}
    After the oracle, the state is:
    \begin{equation*}
        \ket{\psi_2} = \frac{1}{\sqrt{2^n}} \sum_{x \in \left\{0,1\right\}^n} (-1)^{w \cdot x} \ket{x} \ket{-}
    \end{equation*}
    Applying $H^{\otimes n}$ gives the amplitude of each basis state \ket{y}:
    \begin{equation*}
        A(y) = \frac{1}{2^n} \sum_{x \in \left\{0,1\right\}^n} (-1)^{x \cdot (w \oplus y)}
    \end{equation*}
    For $y = w$, the sum equals $2^n$, so the amplitude is $1$. For every $y \neq w$, the positive and negative terms cancel, so the amplitude is $0$. Therefore:
    \begin{equation*}
        \ket{\psi_{\text{final}}} = \ket{w} \ket{-}
    \end{equation*}
    And measurement recovers the hidden string $w$ with probability $1$ (deterministically).
\end{takeawaysbox}